\documentclass[11pt]{article}

\usepackage[a4paper,margin=1in]{geometry}
\usepackage{amsmath,amssymb,amsthm,mathtools}
\usepackage{microtype}
\usepackage[hidelinks]{hyperref}

\allowdisplaybreaks
\numberwithin{equation}{section}

\newtheorem{theorem}{Theorem}[section]
\newtheorem{lemma}[theorem]{Lemma}
\newtheorem{proposition}[theorem]{Proposition}
\newtheorem{corollary}[theorem]{Corollary}
\newtheorem{remark}[theorem]{Remark}

\newcommand{\Z}{\mathbb Z}
\newcommand{\N}{\mathbb N}

\title{An Infinite Pellian Family of Integer Solutions to\\
\texorpdfstring{$y^2+x^2y+z^2x+1=0$}{y^2+x^2 y+z^2 x+1=0}}
\author{}
\date{}

\begin{document}
\maketitle

\begin{abstract}
We give a completely explicit infinite family of integer solutions to
\[
        y^2+x^2y+z^2x+1=0.
\]
The construction fixes $x=-145$ and reduces the equation to a Pell-type conic
\[
        S^2-580Z^2=21025^2-4.
\]
A single integral point on this conic, together with the Pell unit
$289+12\sqrt{580}$ of norm $1$, generates infinitely many further integral
points. A parity argument then converts all of them back to integer solutions
of the original equation.
\end{abstract}

\section{Statement of the result}

The purpose of this note is not to classify all integer solutions. Instead, we
prove the following unconditional statement, which is enough to settle the
``finite versus infinite'' alternative for the given Diophantine equation.

\begin{theorem}\label{thm:main}
The Diophantine equation
\begin{equation}\label{eq:main}
        y^2+x^2y+z^2x+1=0
\end{equation}
has infinitely many integer solutions $(x,y,z)\in\Z^3$.
More precisely, let
\begin{equation}\label{eq:initial-SZ}
        S_0=39551,\qquad Z_0=1391,
\end{equation}
and define sequences $(S_n)_{n\geq 0}$ and $(Z_n)_{n\geq 0}$ recursively by
\begin{equation}\label{eq:recurrence}
\begin{aligned}
        S_{n+1}&=289S_n+6960Z_n,\\
        Z_{n+1}&=12S_n+289Z_n.
\end{aligned}
\end{equation}
Then, for every $n\geq 0$,
\begin{equation}\label{eq:family}
        (x_n,y_n,z_n)
        =\left(-145,\frac{S_n-21025}{2},Z_n\right)
\end{equation}
is an integer solution of \eqref{eq:main}. The solutions obtained in this way
are pairwise distinct. Since \eqref{eq:main} contains $z$ only through $z^2$,
we also get the companion solutions
\begin{equation}\label{eq:family-pm}
        \left(-145,\frac{S_n-21025}{2},-Z_n\right)
        \qquad(n\geq 0).
\end{equation}
\end{theorem}

The first few solutions produced by \eqref{eq:family} are
\begin{equation*}
        (-145,9263,1391),
        \qquad
        (-145,10545287,876611),
        \qquad
        (-145,6101221823,506679767).
\end{equation*}

\section{Reduction to a Pell-type conic}

We first record the exact equivalence between the original equation with
$x=-145$ and a quadratic Pell-type equation.

\begin{lemma}[Completing the square]\label{lem:completion}
Let $z\in\Z$ and $y\in\Z$. Put
\begin{equation}\label{eq:S-def}
        S=2y+21025.
\end{equation}
Then
\begin{equation}\label{eq:fixed-x-equation}
        y^2+21025y-145z^2+1=0
\end{equation}
holds if and only if
\begin{equation}\label{eq:pell-conic}
        S^2-580z^2=21025^2-4.
\end{equation}
Conversely, if $S,z\in\Z$ satisfy \eqref{eq:pell-conic} and $S$ is odd, then
\begin{equation}\label{eq:y-from-S}
        y=\frac{S-21025}{2}
\end{equation}
is an integer satisfying \eqref{eq:fixed-x-equation}.
\end{lemma}

\begin{proof}
Starting from \eqref{eq:S-def}, we have
\begin{equation*}
        S^2=(2y+21025)^2=4y^2+4\cdot 21025\,y+21025^2.
\end{equation*}
Therefore
\begin{align*}
        S^2-580z^2-(21025^2-4)
        &=4y^2+4\cdot 21025\,y+21025^2-580z^2-21025^2+4\\
        &=4\bigl(y^2+21025y-145z^2+1\bigr).
\end{align*}
Hence \eqref{eq:fixed-x-equation} is equivalent to \eqref{eq:pell-conic}.

For the converse, assume $S,z\in\Z$ satisfy \eqref{eq:pell-conic} and that
$S$ is odd. Since $21025$ is also odd, $S-21025$ is even, so
$y=(S-21025)/2$ is an integer. Substituting $S=2y+21025$ into
\eqref{eq:pell-conic} and using the first part of the proof gives
\eqref{eq:fixed-x-equation}.
\end{proof}

Because $x=-145$ gives $x^2=21025$, equation \eqref{eq:fixed-x-equation} is
exactly the original equation \eqref{eq:main} after $x$ has been fixed equal to
$-145$:
\begin{equation*}
        y^2+x^2y+z^2x+1
        =y^2+21025y-145z^2+1.
\end{equation*}
Thus, to prove Theorem \ref{thm:main}, it remains to construct infinitely many
integer pairs $(S,z)$ satisfying \eqref{eq:pell-conic}, with $S$ odd.

\section{A norm-preserving Pell transformation}

Let
\begin{equation}\label{eq:Q-def}
        Q(S,Z)=S^2-580Z^2.
\end{equation}
The Pell-type equation \eqref{eq:pell-conic} is simply
\begin{equation*}
        Q(S,Z)=21025^2-4.
\end{equation*}
We shall use a norm-one automorphism of the binary quadratic form $Q$.

\begin{lemma}[The Pell unit]\label{lem:unit}
The integers
\begin{equation*}
        U=289,
        \qquad
        V=12
\end{equation*}
satisfy
\begin{equation}\label{eq:unit}
        U^2-580V^2=1.
\end{equation}
Equivalently,
\begin{equation*}
        289^2-580\cdot 12^2=1.
\end{equation*}
\end{lemma}

\begin{proof}
This is a direct calculation:
\begin{equation*}
        289^2=83521,
        \qquad
        580\cdot 12^2=580\cdot 144=83520,
\end{equation*}
so $289^2-580\cdot 12^2=83521-83520=1$.
\end{proof}

\begin{lemma}[Norm preservation]\label{lem:norm-preserve}
For every pair $(S,Z)\in\Z^2$, define
\begin{equation}\label{eq:T-def}
        T(S,Z)=(S',Z')=(289S+6960Z,\,12S+289Z).
\end{equation}
Then
\begin{equation}\label{eq:Q-preserved}
        Q(S',Z')=Q(S,Z).
\end{equation}
In particular, if $(S,Z)$ solves
\begin{equation*}
        S^2-580Z^2=21025^2-4,
\end{equation*}
then $T(S,Z)$ also solves the same equation.
\end{lemma}

\begin{proof}
We prove a slightly more general identity. Suppose $U,V$ are integers and put
\begin{equation*}
        S'=US+580VZ,
        \qquad
        Z'=VS+UZ.
\end{equation*}
Then
\begin{align*}
        S'^2-580Z'^2
        &=(US+580VZ)^2-580(VS+UZ)^2\\
        &=U^2S^2+2\cdot 580UVSZ+580^2V^2Z^2\\
        &\qquad{}-580V^2S^2-2\cdot 580UVSZ-580U^2Z^2\\
        &=(U^2-580V^2)S^2-580(U^2-580V^2)Z^2\\
        &=(U^2-580V^2)(S^2-580Z^2).
\end{align*}
Taking $U=289$ and $V=12$, Lemma \ref{lem:unit} gives
$U^2-580V^2=1$. Hence
\begin{equation*}
        S'^2-580Z'^2=S^2-580Z^2.
\end{equation*}
Since $580V=580\cdot 12=6960$, this is exactly the transformation
\eqref{eq:T-def}. The final assertion follows immediately.
\end{proof}

The transformation \eqref{eq:T-def} is the same as multiplication by the unit
$289+12\sqrt{580}$ in the quadratic order $\Z[\sqrt{580}]$:
\begin{equation*}
        S'+Z'\sqrt{580}
        =(S+Z\sqrt{580})(289+12\sqrt{580}).
\end{equation*}
The proof above expands this identity in elementary integer arithmetic, so no
algebraic number theory is being used as a black box.

\section{The initial point and the infinite orbit}

We now verify the initial point on the Pell-type conic and prove that its orbit
under the transformation from Lemma \ref{lem:norm-preserve} gives infinitely
many admissible points.

\begin{lemma}[The initial point]\label{lem:initial-point}
The pair
\begin{equation*}
        (S_0,Z_0)=(39551,1391)
\end{equation*}
satisfies
\begin{equation}\label{eq:initial-solves}
        S_0^2-580Z_0^2=21025^2-4.
\end{equation}
\end{lemma}

\begin{proof}
Again this is direct arithmetic. We compute
\begin{equation*}
        39551^2=1564281601,
        \qquad
        1391^2=1934881,
\end{equation*}
and hence
\begin{equation*}
        580\cdot 1391^2=580\cdot 1934881=1122230980.
\end{equation*}
Therefore
\begin{equation*}
        39551^2-580\cdot 1391^2
        =1564281601-1122230980
        =442050621.
\end{equation*}
On the other hand,
\begin{equation*}
        21025^2=442050625,
\end{equation*}
so
\begin{equation*}
        21025^2-4=442050625-4=442050621.
\end{equation*}
The two values agree, proving \eqref{eq:initial-solves}.
\end{proof}

\begin{proposition}[The recursively generated conic points]\label{prop:conic-points}
Let $(S_n,Z_n)_{n\geq 0}$ be defined by \eqref{eq:initial-SZ} and
\eqref{eq:recurrence}. Then, for every $n\geq 0$,
\begin{equation}\label{eq:all-conic-points}
        S_n^2-580Z_n^2=21025^2-4.
\end{equation}
Moreover, every $S_n$ is odd, and the sequence $(S_n)_{n\geq 0}$ is strictly
increasing. In particular, the pairs $(S_n,Z_n)$ are pairwise distinct.
\end{proposition}

\begin{proof}
First, Lemma \ref{lem:initial-point} proves \eqref{eq:all-conic-points} for
$n=0$. If \eqref{eq:all-conic-points} holds for some $n$, then
\eqref{eq:recurrence} says precisely that
\begin{equation*}
        (S_{n+1},Z_{n+1})=T(S_n,Z_n),
\end{equation*}
where $T$ is the transformation from Lemma \ref{lem:norm-preserve}. Therefore
Lemma \ref{lem:norm-preserve} gives
\begin{equation*}
        S_{n+1}^2-580Z_{n+1}^2=S_n^2-580Z_n^2=21025^2-4.
\end{equation*}
By induction, \eqref{eq:all-conic-points} holds for all $n\geq 0$.

Next we prove the parity claim. Since $S_0=39551$ is odd and
\begin{equation*}
        S_{n+1}=289S_n+6960Z_n,
\end{equation*}
we have, modulo $2$,
\begin{equation*}
        S_{n+1}\equiv S_n \pmod 2,
\end{equation*}
because $289\equiv 1\pmod 2$ and $6960\equiv 0\pmod 2$. Thus all $S_n$ are
odd.

It remains to prove strict increase. We first show that $S_n>0$ and $Z_n>0$ for
all $n\geq 0$. This is true at $n=0$ because $39551>0$ and $1391>0$. If
$S_n>0$ and $Z_n>0$, then the recurrence gives
\begin{equation*}
        S_{n+1}=289S_n+6960Z_n>0,
        \qquad
        Z_{n+1}=12S_n+289Z_n>0.
\end{equation*}
Hence $S_n,Z_n$ are positive for all $n$. Finally,
\begin{equation*}
        S_{n+1}-S_n=288S_n+6960Z_n>0,
\end{equation*}
so $S_{n+1}>S_n$ for every $n\geq 0$. Therefore the sequence $(S_n)$ is
strictly increasing, and the pairs $(S_n,Z_n)$ are pairwise distinct.
\end{proof}

\section{Proof of the main theorem}

We now assemble the preceding lemmas.

\begin{proof}[Proof of Theorem \ref{thm:main}]
For each $n\geq 0$, Proposition \ref{prop:conic-points} gives
\begin{equation*}
        S_n^2-580Z_n^2=21025^2-4,
\end{equation*}
and it also gives that $S_n$ is odd. Hence Lemma \ref{lem:completion} applies
with $z=Z_n$ and
\begin{equation*}
        y_n=\frac{S_n-21025}{2}.
\end{equation*}
Because both $S_n$ and $21025$ are odd, $y_n$ is an integer. Lemma
\ref{lem:completion} then gives
\begin{equation*}
        y_n^2+21025y_n-145Z_n^2+1=0.
\end{equation*}
Since $(-145)^2=21025$, this is the same as
\begin{equation*}
        y_n^2+(-145)^2y_n+Z_n^2(-145)+1=0.
\end{equation*}
Thus
\begin{equation*}
        (x_n,y_n,z_n)=\left(-145,\frac{S_n-21025}{2},Z_n\right)
\end{equation*}
satisfies the original equation \eqref{eq:main}.

By Proposition \ref{prop:conic-points}, the sequence $(S_n)$ is strictly
increasing. Therefore the corresponding integers
\begin{equation*}
        y_n=\frac{S_n-21025}{2}
\end{equation*}
are also strictly increasing. Consequently the triples in \eqref{eq:family} are
pairwise distinct. This gives infinitely many integer solutions.

Finally, equation \eqref{eq:main} depends on $z$ only through $z^2$. Therefore,
whenever $(-145,y_n,Z_n)$ is a solution, so is $(-145,y_n,-Z_n)$. This proves
\eqref{eq:family-pm} as well.
\end{proof}

\section{Direct check of the first solution}

For completeness, we directly verify the first solution without appealing to
the general construction. When $n=0$,
\begin{equation*}
        y_0=\frac{39551-21025}{2}=9263,
        \qquad
        z_0=1391.
\end{equation*}
Then
\begin{align*}
        y_0^2+(-145)^2y_0+z_0^2(-145)+1
        &=9263^2+21025\cdot 9263-145\cdot 1391^2+1\\
        &=85803169+194754575-280557745+1\\
        &=0.
\end{align*}
This agrees with the theorem and illustrates the mechanism of the construction.

\begin{remark}
The proof is entirely constructive. The recurrence \eqref{eq:recurrence}
provides a deterministic algorithm that outputs a new integer solution from the
previous one. No unproved assumption, search result, or classification theorem
is used.
\end{remark}

\end{document}
